\documentclass[11pt,oneside]{book}
\usepackage[english, italian]{babel} 
\usepackage{microtype}
\usepackage[utf8]{inputenc}
\usepackage{subfigure}
\usepackage{epsfig}
\usepackage{amsmath, amssymb}
\usepackage[linesnumbered,lined,commentsnumbered,italiano]{algorithm2e}
\usepackage{fancybox}
\usepackage{listings}
\usepackage{url}
\usepackage{graphicx}
\usepackage{todonotes}
\usepackage{moreverb}
\usepackage{lmodern}
\usepackage{multirow}
\usepackage[labelfont=it,textfont={it}]{caption}
\usepackage[final,bookmarks,breaklinks,colorlinks,allcolors=black]{hyperref}
\usepackage{color}
\usepackage[block=ragged,backend=biber,natbib=true]{biblatex} % Use the bibtex backend with the authoryear citation style (which resembles APA)
\addbibresource{bibliography.bib} % The filename of the bibliography
\usepackage[autostyle=true]{csquotes} % Required to generate language-dependent quotes in the bibliography
\usepackage{hyperref}
\usepackage{enumitem}

%%%%%%%%%%%%%%%%%%%%%%%%%%%%
\definecolor{mygreen}{rgb}{0,0.6,0}
\definecolor{mygray}{rgb}{0.5,0.5,0.5}
\definecolor{mymauve}{rgb}{0.58,0,0.82}
\lstset{ %
	backgroundcolor=\color{white},   % choose the background color; you must add \usepackage{color} or \usepackage{xcolor}; should come as last argument
	basicstyle=\footnotesize,        % the size of the fonts that are used for the code
	breakatwhitespace=false,         % sets if automatic breaks should only happen at whitespace
	breaklines=true,                 % sets automatic line breaking
	captionpos=b,                    % sets the caption-position to bottom
	commentstyle=\color{mygreen},    % comment style
	deletekeywords={...},            % if you want to delete keywords from the given language
	escapeinside={\%*}{*)},          % if you want to add LaTeX within your code
	extendedchars=true,              % lets you use non-ASCII characters; for 8-bits encodings only, does not work with UTF-8
	frame=single,
	keepspaces=true,                 % keeps spaces in text, useful for keeping indentation of code (possibly needs columns=flexible)
	keywordstyle=\color{blue},       % keyword style
	morekeywords={*,...},            % if you want to add more keywords to the set
	numbers=left,                    % where to put the line-numbers; possible values are (none, left, right)
	numberstyle=\tiny\color{mygray}, % the style that is used for the line-numbers
	rulecolor=\color{black},         % if not set, the frame-color may be changed on line-breaks within not-black text (e.g. comments (green here))
	showspaces=false,                % show spaces everywhere adding particular underscores; it overrides 'showstringspaces'
	showstringspaces=false,          % underline spaces within strings only
	showtabs=false,                  % show tabs within strings adding particular underscores
	stringstyle=\color{mymauve},     % string literal style
	tabsize=2,	                   % sets default tabsize to 2 spaces
	title=\lstname                   % show the filename of files included with \lstinputlisting; also try caption instead of title
}
%%%%%%%%%%%%%
\definecolor{gray}{rgb}{0.4,0.4,0.4}
\definecolor{darkblue}{rgb}{0.0,0.0,0.6}
\definecolor{cyan}{rgb}{0.0,0.6,0.6}

%%%%%%%%%%%%%%%%%%%%%%%%
\begin{document}
\selectlanguage{italian}

\begin{titlepage}
\begin{center}
\epsfig{file=Figure/logo_standard.jpg,width=2.5truecm}\\[0.2truecm]
{\Large Universit\`a degli Studi di Salerno}\\[0.2truecm]
{\large Dipartimento di Informatica}\\
\hrulefill
\vfill
{\Large GameUp }\\[0.2truecm]
\vfill\vfill
{\Huge Test Plan Document}
\vfill\vfill


{\bf Studenti} \hfill {\bf Docente}\ \ \\
Francesco Foglia \hfill Prof. Andrea De Lucia\\

\vfill
\hrulefill 

Anno Accademico 2020-2021

\end{center}
\end{titlepage}

\pagenumbering{roman}
\chapter*{Cronologia Revisioni}
\begin{center}
	\begin{tabular}{||c c p{10cm}||} 
	\hline
	Data & Versione & Descrizione \\ [0.5ex] 
	\hline\hline
	26/01/2021 & 1.0 & Release iniziale \\
	\hline
   \end{tabular}
\end{center}

\tableofcontents
\pagestyle{plain}

%%%%%%%%%%%%%%%%%%%%%%%%%%%%
\chapter{Introduzione}
\setcounter{page}{1} 	% devono seguire solo il primo capitolo
\pagenumbering{arabic}	% devono seguire solo il primo capitolo
A causa del tempo ridotto, verranno eseguiti solamente dei test per il rilevamento di errori delle principali funzionalità implementate, ovvero registrazione, login e creazione di una richiesta di pubblicazione. Le ulteriori funzionalità implementate verranno testate nel prossimo processo iterativo di sviluppo, assieme allo sviluppo delle ulteriori funzionalità richieste.

%%%%%%%%%%%%%%%%%%%%%%%%%%%%
\chapter{Approccio}
L'approccio utilizzato sarà di tipo black-box: verranno testate le funzionalità implementate specificate nel RAD senza considerare l'implementazione del codice. Il testing verrà eseguito in modo bottom-up, per evitare l'implementazione di stub ritenuti costosi a causa del paradigma \emph{object-oriented} utilizzato, partendo dalle Repositories che fungono da astrazione per gli oggetti di accesso ai dati permanenti, implementando degli Unit test che testano i singoli metodi in maniera procedurale, poiché tali classi non dipendono da alcuno stato. Proseguendo, verranno verificati i vari Services implementati, senza l'utilizzo di stub poiché sono strettamente correlate alle repositories utilizzate. Terminando, verranno testati gli oggetti Control e le interfacce grafiche (le \emph{view}) che essi ritornano, verificando che l'interfaccia tra la parte front-end e la back-end permetta effettivamente la comunicazione e la chiamata di procedure per effettuare azioni come la visualizzazione o il salvataggio di dati. Gli input per i vari test verranno scelti partizionando l'intero insieme di possibili input in classi di equivalenza, per ridurre il numero di input da utilizzare per la scelta dei vari test cases.

%%%%%%%%%%%%%%%%%%%%%%%%%%%%
\chapter{Test Case Plan}
\section{Registrazione}

\begin{tabular}{|| l | p{25em} ||} 
	\hline
	Parametro & Username\\
	\hline
	Esistente [UE] & \begin{itemize}
	\item L'username esiste già nel database [Error]
	\item L'username è unico [UsernameUnico] 
	\end{itemize}\\
	\hline
\end{tabular}
\\
\begin{tabular}{|| l | p{25em} ||} 
	\hline
	Parametro & Email\\
	\hline
	Esistente [EE] & \begin{itemize}
	\item L'email esiste già nel database [Error]
	\item L'email è unica [EmailUnica, if UsernameUnico] 
	\end{itemize}\\
	\hline
\end{tabular}
\\
\begin{tabular}{|| l | p{25em} ||} 
	\hline
	Parametro & Avatar\\
	\hline
	Presente [PA] & \begin{itemize}
	\item L'avatar è effettivamente una istanza di UploadedFile [AvatarPresente, if UsernameUnico and EmailUnica]
	\item L'avatar non viene passato ed è quindi null [AvatarAssente, if UsernameUnico and EmailUnica]
	\end{itemize}\\
	\hline
\end{tabular}
\\\\\\
\begin{tabular}{|| c | c | c ||}
	\hline
	Codice & Combinazione & Esito\\
	\hline
	TC1\_1 & UE1 & Errore: Username esistente \\
	\hline
	TC1\_2 & UE2,EE1 & Errore: Email esistente \\
	\hline
	TC1\_3 & UE2,EE2,PA1 & Successo: Utente registrato (con avatar) \\
	\hline
	TC1\_4 & UE2,EE2,PA2 & Successo: Utente registrato (ma senza avatar) \\
	\hline
\end{tabular}


\section{Login}
\begin{tabular}{|| l | p{25em} ||} 
	\hline
	Parametro & Username\\
	\hline
	Esistente [UE] & \begin{itemize}
	\item L'username non è presente nel database [Error] 
	\item L'username esiste nel database [UsernameEsiste]
	\end{itemize}\\
	\hline
\end{tabular}
\\
\begin{tabular}{|| l | p{25em} ||} 
	\hline
	Parametro & Password\\
	\hline
	Corretta [PC] & \begin{itemize}
	\item L'hash della password fornita non coincide con l'hash salvato nel database per l'username fornito [Error]
	\item L'hash della password fornita è uguale all'hash associato all'username fornito [PasswordCorretta]
	\end{itemize}\\
	\hline
\end{tabular}
\\\\\\
\begin{tabular}{|| c | c | c ||}
	\hline
	Codice & Combinazione & Esito\\
	\hline
	TC2\_1 & UE1 & Errore: Username non esistente \\
	\hline
	TC1\_2 & UE2,PC1 & Errore: Password non corretta \\
	\hline
	TC1\_3 & UE2,PC2 & Successo: Password corretta, utente autenticato. \\
	\hline
\end{tabular}

\section{Richiesta Pubblicazione Videogioco}
\begin{tabular}{|| l | p{25em} ||} 
	\hline
	Parametro & ID Autore\\
	\hline
	Autore Esistente [AE] & \begin{itemize}
	\item L'ID dell'autore fornito come argomento non è associato ad alcun utente attualmente nel sistema. [Error]
	\item L'ID dell'autore fornito come argomento è effettivamente associato ad un utente attualmente presente nel database. [AE\_OK]
	\end{itemize}\\
	\hline
\end{tabular}
\\\\\\
\begin{tabular}{|| c | c | c ||}
	\hline
	Codice & Combinazione & Esito\\
	\hline
	TC3\_1 & AE1 & Errore: Utente autore non trovato \\
	\hline
	TC3\_2 & AE2 & Successo: Richiesta pubblicazione generata \\
	\hline
\end{tabular}


%%%%%%%%%%%%%%%%%%%%%%%%%%%%
\chapter{Test Case Specification}
\begin{tabular}{|| l | p{25em} ||} 
	\hline
	Identificatore test case & TC1\\
	\hline
	Locazione & tests/Feature/UtenzaRepositoryRegistrazioneTest\\
	\hline
	Feature da testare & La funzionalità da testare è il corretto salvataggio dei dati di un utente nel database relazionale.\\
	\hline
	Criteri di successo/fallimento & Il test passa se viene correttamente inserito nel database un nuovo utente con i dati forniti come input, e in seguito se l'inserimento di un utente con lo stesso username ed email lancia una eccezione.\\
	\hline
	Metodo di controllo & Il test viene avviato tramite il comando \emph{./vendor/sail/bin artisan test}, il quale esegue il TestCase e PHPUnit.\\
	\hline
	Dati & I dati di input vengono forniti tramite un dataProvider che ritorna un iteratore di argomenti che verranno passati alla procedura di test, i quali consistono in due set: uno contenente username, email e password, il secondo contenente username, email, password ed un avatar mock.\\
	\hline
	Procedura di test & Il test viene avviato tramite un comando apposito e la libreria PHPUnit esegue il codice di setup, per poi controllando (tramite l'uso di asserzioni) se i risultati coincidono con l'oracolo proveniente dalle specifiche del requisito in analisi.\\
	\hline
	Requisiti Speciali & I container Docker devono essere in esecuzione, in particolare quello ospitante il database relazionale MySQL. Tutte le migrazioni devono essere state eseguite tramite il comando \emph{./vendor/sail/bin artisan migrate}.\\
	\hline
\end{tabular}
\newpage
\begin{tabular}{|| l | p{25em} ||} 
	\hline
	Identificatore test case & TC2\\
	\hline
	Locazione & tests/Feature/UtenzaRepositoryLoginTest\\
	\hline
	Feature da testare & La funzionalità da testare è la corretta autenticazione di un utente al sistema.\\
	\hline
	Criteri di successo/fallimento & Il test passa se l'utente viene autenticato correttamente all'account associato all'username fornito, se la password associata coincide con quella salvata nel database. Se l'username non è registrato o se la password fornita non è associata all'username fornito, allora viene individuata una failure.\\
	\hline
	Metodo di controllo & Il test viene avviato tramite il comando \emph{./vendor/sail/bin artisan test}, il quale esegue il TestCase e PHPUnit.\\
	\hline
	Dati & I dati di input vengono forniti tramite un dataProvider che ritorna un iteratore di argomenti che verranno passati alla procedura di test, i quali consistono in tre set: uno contenente username non esistente e una password, il secondo contenente l'username dell'account iniziale del sistema ('admin') e una password non corretta ('errore') e il terzo contenente le credenziali corrette per l'account amministratore (username 'admin' e password 'root').\\
	\hline
	Procedura di test & Il test viene avviato tramite un comando apposito e la libreria PHPUnit esegue il codice di setup, per poi controllando (tramite l'uso di asserzioni) se i risultati coincidono con l'oracolo proveniente dalle specifiche del requisito in analisi.\\
	\hline
	Requisiti Speciali & I container Docker devono essere in esecuzione, in particolare quello ospitante il database relazionale MySQL. Tutte le migrazioni devono essere state eseguite tramite il comando \emph{./vendor/sail/bin artisan migrate}. Il \emph{seeding} iniziale del database deve essere effettuato tramite il comando \emph{./vendor/sail/bin artisan db:seed} per permettere la creazione dell'account amministratore.\\
	\hline
\end{tabular}

\newpage
\begin{tabular}{|| l | p{25em} ||} 
	\hline
	Identificatore test case & TC3\\
	\hline
	Locazione & tests/Feature/\newline UtenzaRepositoryCreateRichiestaPubblicazioneTest\\
	\hline
	Feature da testare & La funzionalità da testare è la corretta generazione di una richiesta di pubblicazione di un videogioco.\\
	\hline
	Criteri di successo/fallimento & Il test passa se .\\
	\hline
	Metodo di controllo & Il test viene avviato tramite il comando \emph{./vendor/sail/bin artisan test}, il quale esegue il TestCase e PHPUnit.\\
	\hline
	Dati & I dati di input vengono forniti tramite un dataProvider che ritorna un iteratore di argomenti che verranno passati alla procedura di test, i quali consistono in tre set: uno contenente username non esistente e una password, il secondo contenente l'username dell'account iniziale del sistema ('admin') e una password non corretta ('errore') e il terzo contenente le credenziali corrette per l'account amministratore (username 'admin' e password 'root').\\
	\hline
	Procedura di test & Il test viene avviato tramite un comando apposito e la libreria PHPUnit esegue il codice di setup, per poi controllando (tramite l'uso di asserzioni) se i risultati coincidono con l'oracolo proveniente dalle specifiche del requisito in analisi.\\
	\hline
	Requisiti Speciali & I container Docker devono essere in esecuzione, in particolare quello ospitante il database relazionale MySQL. Tutte le migrazioni devono essere state eseguite tramite il comando \emph{./vendor/sail/bin artisan migrate}. Il \emph{seeding} iniziale del database deve essere effettuato tramite il comando \emph{./vendor/sail/bin artisan db:seed} per permettere la creazione dell'account amministratore.\\
	\hline
\end{tabular}

%%%%%%%%%%%%%%%%%%%%%%%%%%%%
\end{document}
